\section{בליעת פוטון וינון אטום המימן בתורת הפרעות תלויה בזמן}

בסיכום זה נמשיך את הדיון בתורת הפרעות תלויה בזמן (\LR{Time-Dependent Perturbation Theory}) ונתמקד בתהליך של בליעת פוטון (\LR{Photo-absorption}) על ידי מערכת קוונטית, ובפרט בינון (\LR{Photo-ionization}) של אטום המימן.

\subsection{חתך פעולה לבליעת פוטון}

\begin{definitionbox}{חתך פעולה לבליעה}
חתך הפעולה $\sigma_{i\to f}$ לתהליך שבו מערכת קוונטית עוברת ממצב התחלתי $|i\rangle$ למצב סופי $|f\rangle$ בעקבות בליעת פוטון נתון על ידי
\[
\sigma_{i\to f}
=
\frac{4\pi^2}{m^2}\,
\frac{\hbar \alpha}{\omega}
\,
\left|
\langle f | 
e^{i\mathbf{k}_\gamma\cdot\mathbf{x}}
\,\boldsymbol{\varepsilon}\cdot\mathbf{p}
| i \rangle
\right|^2
\,
\delta(E_f - E_i - \hbar\omega),
\]
כאשר $m$ היא מסת האלקטרון, $\alpha$ הוא קבוע המבנה הדק (\LR{Fine Structure Constant}),
$\boldsymbol{\varepsilon}$ הוא וקטור הקיטוב של הפוטון,
$\mathbf{k}_\gamma$ מספר הגל של הפוטון,
ו-$\hbar\omega$ היא אנרגיית הפוטון.
\end{definitionbox}

\begin{notebox}{משמעות פיזיקלית}
חתך הפעולה הוא תכונה פנימית של המערכת הקוונטית. כל התלות בשטף הפוטונים מתקזזת,
והגודל $\sigma$ מתאר שטח אפקטיבי שבו האטום ``מציג את עצמו'' לפוטון.
\end{notebox}

\subsection{ינון אטום המימן וסקלות אורך}

נבחן ינון של אטום המימן:  
המצב ההתחלתי $|i\rangle$ הוא מצב קשור של האלקטרון,
והמצב הסופי $|f\rangle$ הוא מצב רציף (אלקטרון חופשי).

כדי שינון יתרחש נדרש
\[
\hbar\omega \gtrsim E_{\text{bind}},
\]
כאשר אנרגיית הקשר היא מסדר גודל
\[
E_{\text{bind}} \sim \frac{\alpha\hbar c}{a_0},
\]
ו-$a_0$ הוא רדיוס בוהר.

מכאן מתקבל
\[
k_\gamma a_0 \sim \alpha \ll 1,
\]
ולכן אורך הגל של הפוטון גדול בהרבה מרדיוס האטום.

\begin{resultbox}{קירוב הדיפול}
כאשר $k_\gamma a_0 \ll 1$ ניתן לפתח
\[
e^{i\mathbf{k}_\gamma\cdot\mathbf{x}}
\simeq
1 + i\mathbf{k}_\gamma\cdot\mathbf{x} + \cdots,
\]
והאיבר המוביל הוא האיבר הקבוע. זהו קירוב הדיפול החשמלי (\LR{Electric Dipole Approximation}).
\end{resultbox}

\subsection{אופרטור הדיפול והצגה שקולה}

נשתמש בזהות
\[
\mathbf{p} = \frac{m}{i\hbar}[\mathbf{x},H_0],
\]
כאשר
\[
H_0=\frac{\mathbf{p}^2}{2m}+V(r).
\]

\begin{definitionbox}{אופרטור הדיפול}
אופרטור הדיפול החשמלי מוגדר כ
\[
\mathbf{d}=\mathbf{x}
\quad
(\text{בהתעלמות מהמטען }-e).
\]
\end{definitionbox}

שימוש ביחס החילוף נותן
\[
\langle f|\boldsymbol{\varepsilon}\cdot\mathbf{p}|i\rangle
=
\frac{m}{i\hbar}
(E_f-E_i)
\,
\boldsymbol{\varepsilon}\cdot
\langle f|\mathbf{d}|i\rangle.
\]

בעזרת פונקציית הדלתא $E_f-E_i=\hbar\omega$ מתקבל

\begin{resultbox}{נוסחת חתך הפעולה בדיפול}
\[
\sigma_{i\to f}
=
4\pi^2
\,\hbar\alpha\,\omega\,
\left|
\boldsymbol{\varepsilon}\cdot
\langle f|\mathbf{d}|i\rangle
\right|^2
\,
\delta(E_f-E_i-\hbar\omega).
\]
\end{resultbox}

\subsection{כללי ברירה}

נכתוב את מצבי אטום המימן כ-$|n\ell m\sigma\rangle$.

\begin{theorembox}{כללי ברירה לדיפול חשמלי}
אלמנט המטריצה
$\langle n'\ell' m'\sigma'| \mathbf{d} |n\ell m\sigma\rangle$
שונה מאפס רק אם:
\begin{itemize}
\item $\Delta \ell = \pm 1$
\item $\Delta m = 0,\pm 1$ (בהתאם לקיטוב)
\item $\sigma'=\sigma$ (הספין אינו משתנה)
\item הזוגיות מתחלפת.
\end{itemize}
\end{theorembox}

\subsection{ינון באנרגיות גבוהות}

כאשר אנרגיית הפוטון גדולה בהרבה מאנרגיית הקשר,
המצב הסופי של האלקטרון ניתן בקירוב כגל חופשי:
\[
\psi_f(\mathbf{x})
=
\frac{1}{L^{3/2}}
e^{i\mathbf{k}_f\cdot\mathbf{x}}
\chi_\sigma,
\]
כאשר $L^3$ הוא נפח נירמול.

מצב היסוד של המימן הוא
\[
\psi_{1s}(\mathbf{x})
=
\frac{1}{\sqrt{\pi a_0^3}}
e^{-r/a_0}.
\]

חישוב מפורש של אלמנט המטריצה נותן אינטגרל מהצורה
\[
\int d^3x\,
e^{i\mathbf{q}\cdot\mathbf{x}}
e^{-r/a_0}
=
\frac{8\pi a_0^3}{(1+q^2a_0^2)^2}.
\]

\subsection{חתך פעולה דיפרנציאלי}

לאחר חיבור צפיפות המצבים מתקבל חתך פעולה דיפרנציאלי
\[
\frac{d\sigma}{d\Omega}
\propto
(\boldsymbol{\varepsilon}\cdot\hat{\mathbf{p}}_f)^2
\,
\frac{p_f^3}{(p_f^2 + 1/a_0^2)^4}.
\]

\begin{notebox}{התפלגות זוויתית}
האלקטרונים נפלטים בעיקר בכיוון הקיטוב של השדה החשמלי.
בכיוון הניצב לקיטוב ההסתברות מתאפסת.
\end{notebox}

\subsection{התנהגות אסימפטוטית}

באנרגיות גבוהות מתקבל
\[
\sigma(\omega)\sim \omega^{-7/2}.
\]

\begin{resultbox}{משמעות פיזיקלית}
ככל שאנרגיית הפוטון עולה,
חתך הפעולה לינון קטן,
והחומר נעשה שקוף יותר לקרינה עתירת אנרגיה.
\end{resultbox}

\subsection{הערה פיזיקלית כללית}

באנרגיות גבוהות מאוד נכנסות פיזיקות נוספות:
פיזור קומפטון (\LR{Compton Scattering})
ויצירת זוג אלקטרון–פוזיטרון,
אך אלו מעבר להיקף הדיון הנוכחי.

\section{חלקיקים זהים (\LR{Identical Particles})}

\subsection{מבוא ועקרון אי-ההבחנה}

בטבע קיימים סוגים מעטים של חלקיקי יסוד (אלקטרונים, קווארקים, פוטונים וכו'). במערכות רב-גופיות רבות, אנו עוסקים בצוסף של חלקיקים מאותו סוג. בניגוד למכניקה הקלאסית, שבה ניתן "לצבוע" חלקיק ולעקוב אחר מסלולו, במכניקה הקוונטית חלקיקים מאותו סוג הם \textbf{בלתי ניתנים להבחנה} (\LR{Indistinguishable}).

נניח מערכת של $n$ חלקיקים זהים המתוארת על ידי פונקציית גל $\psi(\vec{r}_1, \vec{r}_2, \dots, \vec{r}_n)$. בשל אי-ההבחנה, החלפת מיקומיהם (או דרגות החופש שלהם) של שני חלקיקים $i$ ו-$j$ אינה יכולה לשנות את המציאות הפיזיקלית, כלומר את צפיפות ההסתברות $|\psi|^2$.

\begin{definitionbox}{אופרטור ההחלפה (\LR{Exchange Operator})}
נגדיר את אופרטור ההחלפה $P_{ij}$ כפועל על פונקציית הגל ומחליף בין החלקיק ה-$i$ לחלקיק ה-$j$:
\[ P_{ij} \psi(\dots, \vec{r}_i, \dots, \vec{r}_j, \dots) = \psi(\dots, \vec{r}_j, \dots, \vec{r}_i, \dots) \]
מכיוון שהחלפה כפולה מחזירה את המערכת למצבה המקורי ($P_{ij}^2 = I$), הערכים העצמיים של $P_{ij}$ חייבים להיות $\alpha = \pm 1$.
\end{definitionbox}

\subsection{בוזונים, פרמיונים וקשר לספין}

הטבע מחלק את כל החלקיקים לשתי משפחות על פי התנהגותם תחת החלפה. חלוקה זו נקבעת באופן חד-חד-ערכי על ידי הספין של החלקיק (\LR{Spin-Statistics Theorem}).

\begin{theorembox}{קשר הספין והסטטיסטיקה}
\begin{itemize}
    \item \textbf{בוזונים (\LR{Bosons}):} חלקיקים בעלי ספין שלם ($s=0, 1, 2 \dots$). פונקציית הגל שלהם \textbf{סימטרית} להחלפה:
    \[ P_{ij} \psi = +\psi \]
    \item \textbf{פרמיונים (\LR{Fermions}):} חלקיקים בעלי ספין חצי-שלם ($s=1/2, 3/2 \dots$). פונקציית הגל שלהם \textbf{אנטי-סימטרית} להחלפה:
    \[ P_{ij} \psi = -\psi \]
\end{itemize}
\end{theorembox}

\begin{notebox}{חלקיקים מורכבים}
ניתן להתייחס למערכת קשורה של חלקיקים כאל "חלקיק" אחד. לדוגמה, אטום מימן (פרוטון + אלקטרון, שניהם פרמיונים) מתנהג כבוזון במרחקים גדולים, שכן החלפת שני אטומי מימן שקולה להחלפת שני זוגות של פרמיונים, המייצרת פעמיים סימן מינוס ($(-1)^2 = 1$).
\end{notebox}

\subsection{פורמליזם מתמטי: אופרטורי סימטריזציה}

כדי לבנות פונקציית גל חוקית עבור $n$ חלקיקים זהים מתוך פונקציות של חלקיק יחיד $\phi_a(\vec{r})$, עלינו להפעיל אופרטורים שיבטיחו את הסימטריה הנדרשת.

\begin{definitionbox}{אופרטורי ההטלה $S$ ו-$A$}
עבור $n$ חלקיקים, נגדיר את אופרטור הסימטריזציה ($S$) ואופרטור האנטי-סימטריזציה ($A$):
\[ S = \frac{1}{n!} \sum_{P \in S_n} P, \quad A = \frac{1}{n!} \sum_{P \in S_n} \text{sgn}(P) P \]
כאשר $\text{sgn}(P)$ הוא סימן התמורה ($+1$ לתמורה זוגית, $-1$ לאי-זוגית).
\end{definitionbox}

\begin{examplebox}{שני חלקיקים}
עבור שני חלקיקים במצבים $\phi_a$ ו-$\phi_b$:
\begin{itemize}
    \item \textbf{מצב בוזוני:} $\psi_S \propto \phi_a(1)\phi_b(2) + \phi_b(1)\phi_a(2)$
    \item \textbf{מצב פרמיוני:} $\psi_A \propto \phi_a(1)\phi_b(2) - \phi_b(1)\phi_a(2)$
\end{itemize}
\end{examplebox}

\subsection{עקרון האיסור של פאולי}

אחת התוצאות הדרמטיות ביותר של האנטי-סימטריה של פרמיונים היא חוסר היכולת של שני פרמיונים לאכלס את אותו המצב הקוונטי.

\begin{resultbox}{עקרון האיסור של פאולי (\LR{Pauli Exclusion Principle})}
אם ננסה להושיב שני פרמיונים באותו מצב קוונטי ($\phi_a = \phi_b$), נקבל פונקציית גל מתאפסת זהותית:
\[ \psi_A \propto \phi_a(1)\phi_a(2) - \phi_a(1)\phi_a(2) = 0 \]
מכאן נובע ששני פרמיונים זהים אינם יכולים להימצא באותו מצב קוונטי (כולל ספין).
\end{resultbox}

\begin{center}
\begin{tikzpicture}
    % Illustration of Exchange Symmetry
    \draw[thick, blue, ->] (0,0) node[left] {$\vec{r}_1$} .. controls (1,1) .. (2,0) node[right] {$\vec{r}_2$};
    \draw[thick, red, ->] (2,0) .. controls (1,-1) .. (0,0);
    \node at (1,1.3) {$P_{12}$};
    \node[below] at (1,-1.2) {\LR{Exchange of coordinates}};
    
    % Wavefunction sign
    \node at (5,0.5) {$\psi(\vec{r}_2, \vec{r}_1) = \pm \psi(\vec{r}_1, \vec{r}_2)$};
    \draw [decorate,decoration={brace,amplitude=5pt}] (6.5,-0.2) -- (6.5,1.2) node [black,midway,xshift=1.5cm] {\footnotesize $+ \text{ \LR{Bosons}}$ \\ \footnotesize $- \text{ \LR{Fermions}}$};}
\end{tikzpicture}
\end{center}